% Spelersuitleg bij de maandelijkse jokers (#1003), verteld door Coach Rudy.
%
% Bouwen:  ./scripts/build-uitleg-pdf.sh
% Dat script zet de webp-avatars uit src/features/coach/components/rudi_avatars
% om naar png en roept daarna xelatex aan — LaTeX leest zelf geen webp.
%
% Bewust zonder emoji: XeLaTeX kan de kleuren-emojifont van macOS niet
% insluiten. De kleuren spiegelen de app: --lef (violet) voor het risicospel van
% de spelers, --coach (Rudy's warme amber) en --accent (smaragd) uit
% src/app/index.css. De jokers krijgen bij de bouw hun eigen familie; tot die er
% is volgt dit blad het lef-violet, want een joker hoort naast de lef-tip thuis.
\documentclass[10pt,a4paper]{article}

\usepackage{fontspec}
\usepackage[dutch]{babel}
\usepackage[margin=1.1cm]{geometry}
\usepackage{xcolor}
\usepackage{tcolorbox}
\usepackage{enumitem}
\usepackage{multicol}
\usepackage{parskip}
\setlength{\parskip}{3pt}
\usepackage{tikz}
\usepackage{graphicx}

\tcbuselibrary{skins}

\setmainfont{Avenir Next}
% Iets strakkere regelval: zo past de hele uitleg op één blad. Krapper dan het
% blad van #1004 (0,97), want hier staan er drie jokerkaarten naast elkaar.
\linespread{0.94}

% Kleuren uit src/app/index.css — dezelfde families als in de app zelf.
\definecolor{lef}{HTML}{6F56A8}        % --lef: het risicospel van de spelers
\definecolor{lefsoft}{HTML}{EFEAF8}    % --lef-soft
\definecolor{lefline}{HTML}{D4C9EC}    % --lef-line
\definecolor{coachsoft}{HTML}{FDF0DC}  % --coach-soft
\definecolor{coachline}{HTML}{F0D3A3}  % --coach-line
\definecolor{coachdiep}{HTML}{5A3410}  % --coach-diep
\definecolor{inkt}{HTML}{1A2620}       % --ink
\definecolor{inktzacht}{HTML}{586A60}  % --ink-soft-strong
\definecolor{smaragd}{HTML}{0C8A5F}    % --accent
\definecolor{dorst}{HTML}{9A6010}      % --dorst: de traktatie (#1004)

\color{inkt}
\pagestyle{empty}
\setlist[itemize]{leftmargin=1.2em,itemsep=2.5pt,topsep=3pt}

\newcommand{\kop}[1]{%
  \vspace{0.35em}%
  {\color{lef}\large\bfseries #1}\par\vspace{0.1em}%
}

% Rudy's ronde avatar. De bron is vierkant met zwarte hoeken, dus wegknippen
% tot een cirkel — net als .coach-avatar in de app. #2 is de diameter; de
% straal is de helft daarvan, vandaar \dimexpr en niet ,,2#2'' (dat plakt de
% cijfers aan elkaar tot 21cm).
\newcommand{\rudikop}[2][rudi-portret]{%
  \begin{tikzpicture}
    \clip (0,0) circle (\dimexpr#2/2\relax);
    \node[inner sep=0pt] at (0,0) {\includegraphics[width=#2]{#1}};
  \end{tikzpicture}%
}

% Eén sneer van de coach: avatar links, tekstballon rechts, in het coach-amber.
\newcommand{\rudi}[2][rudi-portret]{%
  \begin{tcolorbox}[colback=coachsoft,colframe=coachline,boxrule=1pt,arc=6pt,
    left=8pt,right=10pt,top=5pt,bottom=5pt]
    \noindent
    \begin{minipage}[c]{1.3cm}\rudikop[#1]{1.15cm}\end{minipage}%
    \begin{minipage}[c]{\dimexpr\linewidth-1.4cm\relax}
      {\color{coachdiep}\itshape #2}
    \end{minipage}
  \end{tcolorbox}%
}

% Eén jokerkaart: naam, ondertitel en wat hij doet, in een eigen kader.
\newcommand{\jokerkaart}[3]{%
  \begin{tcolorbox}[colback=white,colframe=lefline,boxrule=1pt,arc=6pt,
    left=9pt,right=9pt,top=5pt,bottom=5pt]
    {\color{lef}\bfseries #1}\par
    {\color{inktzacht}\footnotesize #2}\par\vspace{2pt}
    #3
  \end{tcolorbox}%
}

\begin{document}

% --- Titelblok -------------------------------------------------------------
\begin{tcolorbox}[colback=lefsoft,colframe=lefline,boxrule=1pt,
  arc=6pt,left=12pt,right=12pt,top=7pt,bottom=7pt]
  \noindent
  \begin{minipage}[c]{\dimexpr\linewidth-2.6cm\relax}
    {\color{lef}\LARGE\bfseries Nieuw: uw geheim wapen}\par\vspace{3pt}
    {\color{inktzacht}\'E\'en joker per maand. U kiest zelf welke, en op welke
    match u hem laat vallen.}
  \end{minipage}%
  \hfill
  \begin{minipage}[c]{2.3cm}
    \hfill\rudikop{2.1cm}
  \end{minipage}
\end{tcolorbox}

\rudi{Elke maand krijgt u van mij \'e\'en kaart in handen. \'E\'en. Niet drie,
niet ,,nog eentje want vorige week ging het niet door''. Sommigen sparen hem tot
de laatste speeldag; de rest speelt hem op een dinsdag tegen iemand die net
leert serveren. Ik noteer allebei.}

% --- Wat is het ------------------------------------------------------------
\kop{Waar gaat dit over?}

Bovenop uw dagelijkse lef-tip krijgt u vanaf nu \textbf{\'e\'en joker per
kalendermaand}. U speelt hem \emph{v\'o\'or} de aftrap uit op een geplande
groepsmatch waarin u zelf meedoet, en hij verandert de inzet of de
omstandigheden van die ene partij. Ongebruikt vervalt hij op de eerste van de
volgende maand --- hij rolt niet door.

% --- De drie kaarten -------------------------------------------------------
\kop{De drie kaarten}

\begin{multicols}{3}
\jokerkaart{Schild}{Deze match telt niet}{Verliest u, dan blijft uw rating
staan. \textbf{Wint u ook.} Uw partner en uw tegenstanders spelen gewoon door
voor hun volle mutatie --- alleen u staat er die partij naast.}

\columnbreak

\jokerkaart{Dubbel of niets}{Uw mutatie maal twee}{Precies de lef-tip, maar
betaald uit uw joker in plaats van uit uw dagtegoed. U kunt dus \'e\'en dag per
maand \textbf{twee keer} dubbel spelen. Winst telt dubbel, verlies ook.}

\columnbreak

\jokerkaart{Wissel van kant}{De tegenstanders draaien om}{Forehand naar
backhand: uw tegenstanders spelen deze match aan de andere kant van het veld.
Alleen in het dubbelspel, en het kost geen enkel ratingpunt --- alleen hun
houvast.}
\end{multicols}

\rudi[rudi-gemeen]{Ja, het schild neemt ook uw winst af. Nee, dat is geen fout.
Een verzekering die alleen uitkeert en nooit iets kost, heet een cadeau --- en
dan speelt iedereen hem elke maand blind op. Nu moet u kiezen, en kiezen is
precies waar de meesten van u niet goed in zijn.}

% --- Hoe ------------------------------------------------------------------
\kop{Zo speelt u hem uit}

\begin{enumerate}[leftmargin=1.4em,itemsep=3pt,topsep=3pt]
  \item \textbf{Plan de match} (of open een ronde die de app zelf inplande).
        Jokers werken alleen op groepsmatches met een starttijd.
  \item \textbf{Tik op de jokertegel} op de wedstrijdkaart, naast de lef-tip:
        daar staat of uw kaart van deze maand nog vrij is.
  \item \textbf{Bedenk u gerust} --- tot de eerste bal kunt u hem weer
        intrekken. Daarna ligt hij vast en is hij verbruikt.
\end{enumerate}

% --- Regels ---------------------------------------------------------------
\kop{De spelregels}

\begin{tcolorbox}[colback=white,colframe=lefline,boxrule=1pt,arc=6pt,
  left=11pt,right=11pt,top=6pt,bottom=6pt]
\begin{itemize}
  \item \textbf{E\'en joker per kalendermaand.} Per speler, over al uw groepen
        samen. Op de eerste van de maand ligt er een verse kaart klaar.
  \item \textbf{Schild en Dubbel of niets pas vanaf tien matches.} Zolang uw
        rating nog inloopt, valt er niets zinnigs te verdubbelen of af te
        schermen.
  \item \textbf{Niemand ziet uw kaart v\'o\'or de aftrap} --- behalve
        \emph{Wissel van kant}, die moet iedereen weten. De rest wordt onthuld
        zodra de match begint, zodat niemand kan meeliften.
  \item \textbf{Niet stapelen.} Zet u Dubbel of niets, dan kan de gewone
        lef-tip er op diezelfde match niet meer bij.
  \item \textbf{Gelijkspel?} Dan telt Dubbel of niets niet mee --- een
        onbesliste partij is geen gewonnen gok. Het schild blijft wel staan.
\end{itemize}
\end{tcolorbox}

% --- Zichtbaarheid --------------------------------------------------------
\kop{Waar ziet u het terug?}

Op de \textbf{wedstrijdkaart} en in de \textbf{feed} draagt de match het merk van
de gespeelde joker, met wie hem uitspeelde en hoe het afliep. Op uw
\textbf{profiel} staat uw voorraad en wat u de vorige maanden hebt uitgespeeld.

\rudi[rudi-gemeen]{En ja, dat blijft staan. Ook die keer dat u dubbel inzette
tegen het sterkste koppel van de club omdat u ,,een goed gevoel'' had. Uw goede
gevoel heeft geen rating. Uw tegenstanders wel.}

\begin{center}
{\color{smaragd}\bfseries \'E\'en kaart, \'e\'en maand, uw keuze ---}
{\color{dorst}\bfseries en zet er meteen een drankje bij.}
\end{center}

\end{document}
